\section{Código abierto contra código cerrado: un equilibrio dinámico saludable}

Karpathy es un partidario de larga data del código abierto. Compara el panorama actual con el ámbito de los sistemas operativos:

\begin{importantbox}{Analogía con Linux}
Linux se ejecuta en más del 60\% de las computadoras. La industria siempre necesita una plataforma abierta común, y en el ámbito de la IA sucede lo mismo. Pero la diferencia clave es que los modelos de IA requieren un capex (gasto de capital) enorme, lo que hace que la competencia en código abierto sea más difícil.
\end{importantbox}

Estado actual:
\begin{itemize}
    \item Los modelos de código cerrado llevan una ventaja de aproximadamente 6 a 8 meses
    \item Los modelos de código abierto ya son suficientemente buenos para la mayoría de los casos de uso de consumo, e incluso pueden ejecutarse localmente
    \item La demanda de inteligencia de frontera se concentrará en proyectos grandes del tipo «trabajo de nivel Premio Nobel» o «reescribir Linux de C a Rust»
\end{itemize}

\begin{warningbox}{Riesgo sistémico de la centralización}
Karpathy expresa su preocupación por la centralización de los modelos de código cerrado: «centralization has a very poor track record». Espera que existan más frontier labs y que el código abierto siga existiendo, aunque vaya unos meses por detrás. «En realidad estamos, casi por accidente, en una posición bastante buena.»
\end{warningbox}

\begin{figure}[H]
\centering
\includegraphics[width=0.82\textwidth]{frame-agent-api.jpg}
\caption{Video aproximadamente en 00:13:16: el giro de una smart home app a un API endpoint es el ejemplo intuitivo de Karpathy de la lógica de distribución de software agent-first}
\end{figure}

\subsection{Resumen de la sección}

El código abierto va por detrás del código cerrado unos seis meses, pero esto resulta ser precisamente un equilibrio saludable: el código abierto cubre la mayoría de los casos de uso, mientras que el código cerrado hace avanzar la frontera. El riesgo sistémico de la centralización requiere que el código abierto actúe como una fuerza de contrapeso.

